\section{Methodology}

This section describes the experimental design, including the factorial structure of independent variables (§3.1), workload profiles and data characteristics (§3.2), dependent variables and measurement procedures (§3.3), statistical analysis plan (§3.4), and ethical considerations (§3.5).

\subsection{Experimental Design: Factorial Structure}

The experiment employs a $3 \times 2 \times 4$ full factorial design with 30 replications per cell, yielding 720 experimental trials. The three factors are:

\begin{enumerate}
    \item \textbf{Partition-Key Design (K):} 3 levels
    \item \textbf{Capacity Mode (C):} 2 levels  
    \item \textbf{Workload Profile (W):} 4 levels
\end{enumerate}

\subsubsection{Factor 1: Partition-Key Design (K)}

Three partition-key designs are evaluated, operationalising patterns documented in AWS guidance \citep{aws2025sharding} and prior sharding research \citep{eldeeb2022neuroshard}:

\begin{itemize}
    \item \textbf{K1 (Simple):} Partition key = \texttt{orderId} (UUID v4, high cardinality), no sort key. This design maximises distribution uniformity but provides no secondary access path.
    \item \textbf{K2 (Composite):} Partition key = \texttt{customerId} (10,000 unique values), sort key = \texttt{orderTimestamp} (ISO 8601 timestamp with microsecond precision). This design enables customer-localised range queries (e.g., ``retrieve all orders for customer X in date range Y–Z'') but concentrates orders from high-activity customers on the same partition.
    \item \textbf{K3 (Write-Sharded):} Partition key = \texttt{orderId\#shard}, where shard $\in \{0, 1, \ldots, 9\}$ is appended via modulo hashing (\texttt{shard = hash(orderId) mod 10}). Writes are uniformly distributed across 10 shard partitions. Reads must use a scatter-gather pattern: query all 10 shards with \texttt{orderId\#0}, \texttt{orderId\#1}, \ldots, \texttt{orderId\#9} and merge results client-side. This design trades read amplification (10× read operations per logical read) for write distribution.
\end{itemize}

\subsubsection{Factor 2: Capacity Mode (C)}

Two capacity modes are evaluated:

\begin{itemize}
    \item \textbf{C1 (On-Demand):} \texttt{BillingMode = PAY\_PER\_REQUEST}. Automatic scaling to any traffic level. Billed per request: USD 1.25 per million write request units (WRU), USD 0.25 per million read request units (RRU), as of January 2025 in \texttt{eu-west-1}.
    \item \textbf{C2 (Provisioned):} \texttt{BillingMode = PROVISIONED} with target-tracking auto-scaling. Reserved capacity is provisioned per key design based on anticipated load (see Table \ref{tab:capacity-plan}). Billed per reserved capacity-hour: USD 0.00065 per WCU-hour, USD 0.00013 per RCU-hour, plus consumed write/read units during burst periods exceeding reserved capacity.
\end{itemize}

Table \ref{tab:capacity-plan} shows the provisioned capacity configuration for each key design, derived from design checks (§4.3):

\begin{table}[h]
\centering
\caption{Provisioned Capacity Configuration by Key Design}
\label{tab:capacity-plan}
\begin{tabular}{|l|c|c|c|}
\hline
\textbf{Configuration} & \textbf{Min RCU} & \textbf{Target RCU} & \textbf{Max RCU} \\ \hline
K1 (Simple) & 140 & 210 & 420 \\ \hline
K2 (Composite) & 140 & 210 & 420 \\ \hline
K3 (Write-Sharded) & 1360 & 2040 & 4080 \\ \hline
& \textbf{Min WCU} & \textbf{Target WCU} & \textbf{Max WCU} \\ \hline
K1 (Simple) & 140 & 210 & 420 \\ \hline
K2 (Composite) & 140 & 210 & 420 \\ \hline
K3 (Write-Sharded) & 140 & 210 & 420 \\ \hline
\end{tabular}
\end{table}

K3's elevated read capacity (1360 vs. 140 RCU) accounts for read amplification: each logical read requires querying 10 shard partitions.

\subsubsection{Factor 3: Workload Profile (W)}

Four workload profiles are evaluated, extending the read/write mix patterns of \citet{pantelic2026benchmarking} with a burst profile:

\begin{itemize}
    \item \textbf{W1 (Read-Heavy):} 95\% GetItem, 5\% PutItem, sustained 200 operations/second for 300 seconds (60,000 operations total).
    \item \textbf{W2 (Write-Heavy):} 30\% GetItem, 70\% PutItem, sustained 200 operations/second for 300 seconds.
    \item \textbf{W3 (Mixed):} 50\% GetItem, 50\% PutItem, sustained 200 operations/second for 300 seconds.
    \item \textbf{W4 (Burst):} 50\% GetItem, 50\% PutItem, alternating 60 seconds at 200 ops/s and 30 seconds at 1000 ops/s, repeated for 300 seconds total. This profile stresses adaptive capacity and auto-scaling response time.
\end{itemize}

All workloads are driven by AWS Lambda functions with 1024 MB memory allocation, invoked concurrently to achieve target aggregate throughput. Lambda functions are pre-warmed (5 warm-up invocations per function) before measurement begins to eliminate cold-start latency.

\subsection{Data Characteristics and Access Distribution}

\subsubsection{Dataset}

The dataset comprises 1,000,000 synthetic order records distributed across 10,000 unique customer IDs. Each order record contains:

\begin{itemize}
    \item \texttt{orderId}: UUID v4 (36 characters)
    \item \texttt{customerId}: Integer in range [1, 10000]
    \item \texttt{orderTimestamp}: ISO 8601 timestamp with microsecond precision
    \item \texttt{itemCount}: Integer in range [1, 10]
    \item \texttt{totalAmount}: Decimal (2 decimal places)
    \item \texttt{status}: String (\texttt{pending}, \texttt{confirmed}, \texttt{shipped}, \texttt{delivered})
    \item \texttt{payload}: Random alphanumeric string padded to achieve 1 KB total item size
\end{itemize}

The 1 KB item size is representative of order/transaction records in e-commerce and financial services applications. Item size directly affects consumed capacity: a 1 KB item consumes 1 RCU for eventually consistent reads, 1 WCU for writes.

\subsubsection{Access Distribution}

Operations follow a Zipfian distribution with shape parameter $s = 1.070916$, calibrated such that approximately 90\% of operations access approximately 10\% of keys. This skew models realistic access patterns in which popular items (e.g., trending products, high-volume customers) receive disproportionate traffic \citep{kanellis2025f2, ziegler2023oltp}.

The Zipfian probability mass function is:

\begin{equation}
P(k; s, N) = \frac{k^{-s}}{\sum_{n=1}^{N} n^{-s}}
\end{equation}

where $k$ is the rank of the key, $s$ is the shape parameter, and $N$ is the number of unique keys. The denominator is the generalised harmonic number $H_{N,s}$.

For K1 and K3, the Zipfian distribution is applied over the 1,000,000 unique order IDs. For K2, the distribution is applied over the 10,000 unique customer IDs, and within each customer's orders, a secondary Zipfian distribution (s = 0.8) selects specific orders by timestamp rank. This nested distribution models scenarios in which both customers and orders within customers follow power-law access patterns.

\subsection{Dependent Variables and Measurement}

Six dependent variables are measured for each experimental trial:

\begin{enumerate}
    \item \textbf{Mean Latency (ms):} Arithmetic mean of end-to-end client-observed latency (Lambda invocation to DynamoDB response receipt) for all operations in the batch.
    \item \textbf{p95 Latency (ms):} 95th percentile latency.
    \item \textbf{p99 Latency (ms):} 99th percentile latency. Tail latency is a critical metric for interactive applications \citep{ziegler2023oltp}.
    \item \textbf{Throughput (ops/s):} Successfully completed operations divided by wall-clock time. Excludes throttled operations (which return \texttt{ProvisionedThroughputExceededException} or \texttt{RequestLimitExceeded}).
    \item \textbf{Throttle Rate (\%):} Percentage of operations rejected by DynamoDB admission control. Throttled operations are not retried (Lambda handler configured with \texttt{total\_max\_attempts=1}).
    \item \textbf{Cost per 10,000 Operations (USD):} Computed from consumed RCU/WCU (on-demand mode) or reserved capacity-hours plus burst consumption (provisioned mode). Cost model details are provided in §4.4.
\end{enumerate}

Latency and throttle rate are instrumented in the Lambda handler using the AWS SDK for Python (boto3) with millisecond-resolution timestamps. Consumed capacity is extracted from DynamoDB API responses (\texttt{ConsumedCapacity} field). CloudWatch metrics (\texttt{ConsumedReadCapacityUnits}, \texttt{ConsumedWriteCapacityUnits}, \texttt{UserErrors}) provide secondary validation.

\subsection{Statistical Analysis Plan}

The primary analysis employs two-way analysis of variance (ANOVA) with interaction to test the effects of key design (K), capacity mode (C), and their interaction (K × C) on each dependent variable, separately for each workload profile (W). The model is:

\begin{equation}
Y_{ijkl} = \mu + \alpha_i + \beta_j + (\alpha\beta)_{ij} + \epsilon_{ijkl}
\end{equation}

where:
\begin{itemize}
    \item $Y_{ijkl}$ is the observed value of the dependent variable for key design $i$, capacity mode $j$, workload $k$, replication $l$
    \item $\mu$ is the grand mean
    \item $\alpha_i$ is the main effect of key design (K)
    \item $\beta_j$ is the main effect of capacity mode (C)
    \item $(\alpha\beta)_{ij}$ is the interaction effect
    \item $\epsilon_{ijkl}$ is the residual error, assumed $\sim N(0, \sigma^2)$
\end{itemize}

If ANOVA assumptions (normality, homoscedasticity) are violated (assessed via Shapiro-Wilk test and Levene's test), the aligned rank transform (ART) procedure \citep{wobbrock2011aligned} is applied to convert data to ranks while preserving main effects and interactions.

Post-hoc pairwise comparisons are conducted using Tukey's Honest Significant Difference (HSD) test to identify which specific configurations differ significantly. The family-wise error rate is controlled at $\alpha = 0.05$ using the Holm-Bonferroni correction across 12 pairwise comparisons (6 configurations × (6-1)/2).

\subsection{Randomisation and Replication}

To control for time-of-day effects (DynamoDB multi-tenant load variations, network latency variations), the 24 experimental cells (6 configurations × 4 workloads) are executed in randomised blocks. Each block contains one trial of each cell, executed in random order. 30 blocks are completed, yielding 30 replications per cell. Each trial is separated by a 120-second settling period to allow auto-scaling to stabilise and CloudWatch metrics to propagate.

\subsection{Ethical Considerations}

This study uses synthetic data only; no personal data or personally identifiable information (PII) is collected, stored, or processed. All experiments are conducted on the researcher's personal AWS account under the AWS Free Tier and educational credits, with cost monitoring and budget alerts configured to prevent unintended expenditure. The study complies with National College of Ireland's research ethics policy for non-human-subjects research. No IRB approval is required.

\subsection{Reproducibility and Artefact Availability}

The complete experimental artefact is provided as open-source Infrastructure-as-Code (Terraform), Lambda function handlers (Python 3.12), workload generators, metrics collectors, statistical analysis scripts (Python with SciPy, Statsmodels, Pandas), and configuration management (YAML). Third parties can replicate the experiment by:

\begin{enumerate}
    \item Deploying the Terraform infrastructure to their AWS account (6 DynamoDB tables, 1 Lambda function, IAM roles, CloudWatch alarms).
    \item Seeding the 1,000,000-item dataset using the provided batch-write script.
    \item Executing the randomised block matrix using the provided orchestration script.
    \item Collecting CloudWatch metrics and Lambda-recorded latency/throttle data.
    \item Re-running the statistical analysis scripts to reproduce ANOVA, post-hoc tests, and visualisations.
\end{enumerate}

Configuration parameters (region, pricing, Zipfian shape, replication count) are externalised to \texttt{config/experiment.yaml} to facilitate sensitivity analysis. The artefact is documented in a separate configuration manual (see Appendix).
